\begin{textsample}{NEG Dim 6 – human – Score -3.00 – mg/mg\_ef\_ch\_his\_7\_p2.txt}  \label{ex:f6_neg_003}
A proposta de avaliação considera as habilidades a serem desenvolvidas em cada ano de escolaridade.
O desenvolvimento do raciocínio histórico, da perspectiva temporal e da investigação servirão de parâmetro para a avaliação do desenvolvimento cognitivo dos estudantes.
No entanto, é necessário que o professor esteja atento ao fato de que muitas das capacidades requeridas para o desenvolvimento do raciocínio histórico e da cidadania só serão consolidadas no decorrer de um período maior, \textbf{exigindo} persistência no trabalho com um núcleo comum de habilidades e atitudes por meio de estratégias de ensino e de avaliação, que estabeleçam diferentes graus de complexidade ao longo do Ensino Fundamental.

% matched lemmas: exigir
\end{textsample}
